\chapter{Raw TE Analysis}

\section{Workflow Goal}

The raw TE workflow converts paired ZEM/LFA measurements into processed
transport-property tables. The standard entry point is \code{python main.py analyze}.

\section{Command}
The first step is to run a scrip (\code{scripts/sync_lab_metadata.py}) to sync sample information into a notebook
 (\code{data/lab/lab_metadata.md}), which record detailed information of
experiemnts samples (batch id, comptosition, density, heat capacity, doping elements, etc.). This script will
 automatically detect the new added samples in \code{\data\raw}, and add new template in markdown. The $density$ and 
 $heat\ \ capacity$ must be filled before running TE Analysis program. 

\begin{lstlisting}[language=bash,caption={Sync batch information.}]
  # TODO: sync batch/sample information to \code{data/lab/lab_metadata.md}.
  python scripts/sync_lab_metadata.py
  \end{lstlisting}

\begin{lstlisting}[language=bash,caption={Raw TE analysis command slot.}]
# TODO: replace with the exact batch or sample selector.
python main.py analyze CHY-1048

# Optional selected-sample analysis.
python main.py analyze CHY-1054 --sample B C
\end{lstlisting}

\section{Analysis Output}

\begin{lstlisting}[caption={Raw TE analysis terminal output slot.}]
# TODO: paste the key processing lines here.
# Example:
# step1. Data processing
# Saved processed data to: data/processed/...
# Saved extracted features to: results/...
\end{lstlisting}

Processed data will be located at \code{data/processed}, with extracted raw data, and lattice thermal conductivity ($\kappa_L$), 
quality factor based on SPB model, and ZT value.

\section{Noitons}

\begin{itemize}[leftmargin=*]
  \item Required ZEM/LFA file naming and raw folder structure.
  \item Required metadata: density, \code{cp_value}, sample ID, composition.
  \item Output CSV columns and extracted feature JSON.
  \item Common failure modes: missing LFA/ZEM files, missing density, unit
  mismatch, malformed instrument export.
\end{itemize}
